% -----------------------------------------------------------------------
\section{Introduction}
\label{sec:intro}
% -----------------------------------------------------------------------

Urban airflow modelling is a cornerstone of applications ranging from air-quality
monitoring and emergency dispersion prediction to urban heat-island mitigation
and wind-comfort assessment. Accurate predictions require resolving the
complex turbulent dynamics induced by densely packed building geometries, a task
for which LES at metre-scale resolution is the accepted gold standard. The \palm{}
model system \cite{maronga2020}, for example, resolves boundary-layer flow with
sub-grid-scale closures appropriate for urban canopies. While the physical
fidelity of such solvers is well established, their computational cost---hours of
wall time per simulation on modern HPC hardware---makes them incompatible with
the rapid update cycles demanded by operational data assimilation, where an
ensemble of tens to hundreds of model integrations must be completed per
assimilation window.

Neural surrogates offer a promising alternative. By training a learned model on
an offline library of high-fidelity CFD trajectories, one can amortise the
simulation cost into a one-time training phase and achieve forward-model
evaluations that are orders of magnitude faster at inference time. The challenge
is to build surrogates that generalise faithfully over a physically diverse
ensemble of inflow conditions, preserve mass-conservation constraints
approximately, and remain stable under extended autonomous rollout---conditions
that arise naturally in the data-assimilation setting \cite{brandstetter2022}.

In this work we develop and evaluate a neural surrogate for the Barcelona
urban-block geometry, targeting its use inside an \esmda{} framework. The
surrogate is based on the \pthreed{} architecture \cite{holzschuh2025p3d}, a hybrid
convolutional-attention model that couples a U-Net encoder--decoder with a
windowed self-attention bottleneck, providing both local feature extraction
and long-range spatial context at tractable cost. The model is conditioned on
explicit inflow parameters (wind angle and velocity magnitude)
and trained with a pushforward rollout curriculum that progressively increases the
length of the predicted trajectory during training to reduce exposure to
distribution shift during autonomous rollout. The trained surrogate is embedded in
an \esmda{} loop that jointly updates the velocity-field state and inflow parameters
conditioned on sparse point observations from virtual street-level sensors.
